\documentclass[conference, a4paper]{IEEEtran}
\usepackage[utf8]{inputenc}
\usepackage[spanish]{babel}
\usepackage{amsmath, amssymb, amsfonts, amsthm}
\usepackage{algorithmic}
\usepackage{algorithm}
\usepackage{graphicx}
\usepackage{textcomp}
\usepackage{xcolor}
\usepackage{booktabs}
\usepackage{multirow}
\usepackage{hyperref}
\usepackage{tikz}
\usetikzlibrary{babel}

\begin{document}

% Portada Académica Clásica (1 columna, en página propia)
\begin{titlepage}
    \centering
    \vspace*{2cm}
    
    \includegraphics[width=0.4\textwidth]{logo_univalle.png}\par\vspace{1.5cm}
    
    {\Large\bfseries UNIVERSIDAD DEL VALLE \par}
    {\large\bfseries ESCUELA DE INGENIERÍA DE SISTEMAS Y COMPUTACIÓN \par}
    \vspace{1.5cm}
    
    {\huge\bfseries Calculando el Plan de Riego Óptimo de una Finca \par}
    \vspace{0.5cm}
    
    \vspace{2cm}
    
    {\Large \textbf{Autores:}\par}
    {\Large Ervin Caravali Ibarra \par}
    {\large \texttt{ervin.caravali@correounivalle.edu.co}\par}
    \vspace{0.5cm}
    {\Large Camilo Urrea \par}
    {\large \texttt{camilo.urrea@correounivalle.edu.co}\par}
    
    \vspace{2cm}
    
    {\Large \textbf{Asignatura:}\par}
    {\Large Análisis de Algoritmos II \par}
    
    \vfill
    
    {\large Cali, Colombia\par}
    {\large 2026 \par}
\end{titlepage}

% Índice en 1 columna, en página propia
\onecolumn
\begingroup
  \small
  \tableofcontents
\endgroup
\newpage

% De aquí en adelante a 2 columnas
\twocolumn

\begin{abstract}
Este trabajo compara tres estrategias para calcular el orden de riego que minimice las pérdidas por estrés hídrico en una finca cañera: Fuerza Bruta, un algoritmo Voraz con lookahead predictivo y Programación Dinámica con memoización. Mostramos que el problema pertenece a la clase NP-Hard, lo que hace inviable la Fuerza Bruta para instancias grandes. La Programación Dinámica reduce la complejidad de $O(n!)$ a $O(n \cdot 2^n)$ y produce resultados exactos hasta $n \approx 18$, mientras el Voraz escala sin problemas hasta $n=25$ con errores menores al 3\%. Los tres algoritmos se validaron contra 19 casos de prueba con salidas esperadas conocidas. También se discute cómo integrar este motor algorítmico en una arquitectura web real con frontend en React/Angular, backend en Node.js y lógica de optimización en Python.
\end{abstract}

\renewcommand{\IEEEkeywordsname}{Palabras clave}
\begin{IEEEkeywords}
Programación Dinámica, Algoritmos Voraces, Optimización Combinatoria, Scheduling, Riego Agrícola, NP-Hard.
\end{IEEEkeywords}

\section{Introducción y Contexto}
\subsection{El Problema Real y la Importancia del Riego Óptimo}
La agricultura moderna enfrenta el desafío de maximizar la producción mientras reduce el consumo de recursos finitos, particularmente el agua. En regiones de alta intensidad agrícola como el Valle del Cauca, la gestión eficiente del recurso hídrico no es solo una cuestión ambiental, sino que tiene un impacto directo en la rentabilidad de la finca.

Una finca típica está subdividida en parcelas denominadas \textit{tablones}. En este estudio, el cultivo es estrictamente caña de azúcar. Aunque se trata de un monocultivo, cada tablón alberga cepas en diferentes estadios de desarrollo (plantilla, soca de distintos cortes), lo que implica que cada unidad tiene necesidades hídricas y tolerancias al estrés hídrico distintas. El principal cuello de botella no suele ser la disponibilidad absoluta de agua, sino la infraestructura de distribución: la finca dispone de un único sistema de riego centralizado (como un pivote central o una cuadrilla de irrigación). Esta restricción física impone que los tablones deben atenderse de forma secuencial.

\subsection{Restricciones Prácticas y Estrés Hídrico}
El modelo agrónomo establece que la caña de azúcar tiene una ''ventana ideal'' para recibir irrigación. Iniciar el riego en este \textit{tiempo óptimo} asegura el máximo rendimiento de biomasa y sacarosa. Desviarse de este punto, ya sea por adelanto o retraso, genera penalizaciones fisiológicas para el cultivo.

Si el retraso supera un límite denominado \textit{tiempo de supervivencia}, el tablón entra en déficit hídrico agudo. En este punto la caña comienza a marchitarse y el costo penal, que representa las pérdidas económicas directas, crece en función de la prioridad comercial del tablón.

\subsection{Motivación Tecnológica e Impacto}
Históricamente, los ingenieros agrónomos han dependido de heurísticas intuitivas o programación manual para coordinar estos sistemas. Estas aproximaciones rara vez alcanzan el óptimo global, generando desperdicios económicos y un uso subóptimo del agua. La motivación de este proyecto es dotar al sector cañicultor de herramientas algorítmicas que calculen el orden de riego que minimice la penalización total en la cosecha.

\section{Formalización Matemática}
\subsection{Definición del Espacio del Problema: La Finca}
Sea $F$ el conjunto ordenado finito que representa la finca, compuesto por $n$ tablones indexados de 0 a $n-1$:
\begin{equation}
F = \langle T_0, T_1, \dots, T_{n-1} \rangle
\end{equation}

Cada tablón $T_i$ se representa mediante una tupla con cuatro parámetros:
\begin{equation}
T_i = \langle ts_i, tr_i, p_i, ro_i \rangle
\end{equation}

\begin{itemize}
    \item $ts_i \in \mathbb{N}$: \textbf{Tiempo de supervivencia}. Límite antes de empezar a perder valor.
    \item $tr_i \in \mathbb{N}^+$: \textbf{Tiempo de regado}. Duración de la tarea de riego.
    \item $p_i \in \{1, 2, 3, 4\}$: \textbf{Prioridad del tablón}. Factor de ponderación económica donde 4 es la mayor criticidad.
    \item $ro_i \in \mathbb{N}$: \textbf{Tiempo de riego óptimo}. El día ideal en el cual la penalidad es mínima.
\end{itemize}

Imponemos la restricción:
\begin{equation}
0 \leq ro_i \leq ts_i - tr_i \quad \forall i \in \{0, \dots, n-1\}
\end{equation}

\subsection{La Programación de Riego (Permutación)}
Dado que el riego es secuencial, una solución al problema es una \textbf{permutación} $\Pi$ del conjunto de índices $I = \{0, 1, \dots, n-1\}$.

\begin{equation}
\Pi = \langle \pi_0, \pi_1, \dots, \pi_{n-1} \rangle
\end{equation}

Donde $\pi_k = i$ indica que el tablón $T_i$ ocupará la $k$-ésima posición en la secuencia de riego.

\subsection{Cálculo de Tiempos y Función de Costo Individual (CRF)}
Para la primera posición, el tiempo de inicio es $t_{\pi_0}^\Pi = 0$.
Para cualquier tablón subsiguiente $\pi_j$, su tiempo de inicio es el tiempo de finalización de su predecesor:
\begin{equation}
t_{\pi_j}^\Pi = t_{\pi_{j-1}}^\Pi + tr_{\pi_{j-1}}
\end{equation}

La penalidad de un tablón $i$ depende de $t_i^\Pi$. Se evalúan tres casos mutuamente excluyentes:

\textbf{Caso 1 - Riego en Tiempo Óptimo:}
\begin{equation}
CRF_{\Pi}[i] = ts_i - (t_i^\Pi + tr_i) \quad \text{si } t_i^\Pi = ro_i
\end{equation}

\textbf{Caso 2 - Retraso Tolerado sin Déficit Hídrico:}
\begin{equation}
CRF_{\Pi}[i] = 2 \cdot \left( ts_i - (t_i^\Pi + tr_i) \right)
\end{equation}

\textbf{Caso 3 - Déficit Hídrico Crítico (Peor Escenario):}
\begin{equation}
CRF_{\Pi}[i] = 2 \cdot p_i \cdot \left( (t_i^\Pi + tr_i) - ts_i \right)
\end{equation}

\subsection{Función de Costo Objetivo}
El problema de optimización consiste en encontrar la permutación $\Pi^*$ que minimice el costo total:
\begin{equation}
\Pi^* = \arg\min_{\Pi \in S_n} \left( \sum_{i=0}^{n-1} CRF_{\Pi}[i] \right)
\end{equation}

\section{Naturaleza Combinatoria del Problema}
El problema de encontrar la secuencia de riego óptima es un problema de secuenciamiento en una sola máquina \cite{graham1979optimization, pinedo2016scheduling}. Dado que la función de penalidad es escalonada y depende de la posición en la secuencia, no tiene monotonía.

El dominio de soluciones equivale a todas las ordenaciones de los $n$ tablones, de modo que $|S_n| = n!$.

\begin{table}[htbp]
\centering
\caption{Crecimiento de Permutaciones vs Tamaño ($n$)}
\begin{tabular}{@{}c c c@{}}
\toprule
\textbf{Tamaño ($n$)} & \textbf{Permutaciones ($n!$)} & \textbf{Tiempo aprox. a 1 $\mu s$} \\ \midrule
5 & 120 & 0.12 milisegundos \\
10 & 3.628.800 & 3.6 segundos \\
12 & $4.79 \times 10^8$ & $\sim 8$ minutos \\
15 & $1.30 \times 10^{12}$ & $\sim 15$ días \\
20 & $2.43 \times 10^{18}$ & $\sim 77.000$ años \\ \bottomrule
\end{tabular}
\end{table}

Esta explosión factorial es la razón principal por la que el problema se clasifica como \textbf{NP-Hard} \cite{cormen2009introduction}, y motiva el desarrollo de las estrategias algorítmicas que se presentan a continuación.

\section{Análisis Completo de Ejemplos y Funcionamiento Interno}
Para ilustrar cómo opera cada algoritmo, analizamos el archivo \texttt{test1\_in.txt} del proyecto: 
$n=5$ con $T_0=\langle 10, 10, 1, 0 \rangle$, $T_1=\langle 11, 2, 1, 5 \rangle$, $T_2=\langle 13, 5, 1, 7 \rangle$, $T_3=\langle 9, 4, 4, 1 \rangle$ y $T_4=\langle 14, 1, 2, 10 \rangle$.

Al ejecutar la implementación en Python, \textbf{los tres algoritmos coinciden en la misma solución}:
Permutación $\Pi^* = \langle 2, 3, 1, 4, 0 \rangle$ con un costo mínimo de $\mathbf{44}$.
A continuación se detalla cómo cada enfoque llega a este resultado.

\subsection{Fuerza Bruta: El Árbol de Búsqueda (DFS)}
La Fuerza Bruta genera un árbol recursivo explorando las $n! = 120$ combinaciones posibles. En la \textbf{Figura \ref{fig:fb_tree}}, los caminos en rojo corresponden a decisiones subóptimas, mientras que el camino verde muestra la rama que alcanza el costo mínimo de 44.

El cálculo de la rama óptima es el siguiente:
\begin{itemize}
    \item \textbf{Elige $T_2$ ($t=0$):} Caso 2. $C = 2 \times (13 - (0+5)) = 16$. Nuevo $t=5$. Acumulado = $16$.
    \item \textbf{Elige $T_3$ ($t=5$):} Caso 2. $C = 2 \times (9 - 9) = 0$. Nuevo $t=9$. Acumulado = $16$.
    \item \textbf{Elige $T_1$ ($t=9$):} Caso 2. $C = 2 \times (11 - 11) = 0$. Nuevo $t=11$. Acumulado = $16$.
    \item \textbf{Elige $T_4$ ($t=11$):} Caso 2. $C = 2 \times (14 - 12) = 4$. Nuevo $t=12$. Acumulado = $20$.
    \item \textbf{Elige $T_0$ ($t=12$):} Marchitado (Caso 3). $C = 2 \times 1 \times (22 - 10) = 24$. Total = $\mathbf{44}$.
\end{itemize}

\begin{figure*}[htbp]
\centering
\begin{tikzpicture}[
  scale=0.75, transform shape,
  level distance=1.8cm,
  level 1/.style={sibling distance=3.2cm},
  level 2/.style={sibling distance=2cm},
  every node/.style={rectangle, draw, align=center, fill=white, scale=0.85, font=\bfseries, rounded corners},
  opt/.style={draw=green!60!black, fill=green!10, thick},
  bad/.style={draw=red!80, fill=red!10},
  edge from parent/.style={draw, thick, ->, >=stealth}
]
  \node {Raíz\\($t=0$, Costo=$0$)}
    child {node[opt] {$T_2$\\($t=5$, C=$16$)}
      child {node[opt] {$T_3$\\($t=9$, C=$16$)}
        child {node[opt] {$T_1$\\($t=11$, C=$16$)}
            child {node[opt] {$T_4$\\($t=12$, C=$20$)}
                child {node[opt] {$T_0$\\($t=22$, C=$\mathbf{44}$)}}
            }
        }
        child {node[bad] {$T_0$\\($t=19$, C=$34$)}
            child {node[bad] {\dots\\(Costo $\ge 50$)}}
        }
      }
      child {node[bad] {$T_4$\\($t=6$, C=$24$)}
        child {node[bad] {\dots\\(Costo $\ge 60$)}}
      }
    }
    child {node[bad] {$T_1$\\($t=2$, C=$8$)}
      child {node[bad] {\dots\\(Costo $\ge 70$)}}
    }
    child {node[bad] {$T_0$\\($t=10$, C=$0$)}
      child {node[bad] {\dots\\(Costo $\ge 80$)}}
    };
\end{tikzpicture}
\caption{Visualización parcial del árbol de ejecución (Fuerza Bruta). El camino verde representa la decisión óptima $\Pi^*$ y detalla el tiempo $t$ y el costo acumulado ($C$) en cada nodo.}
\label{fig:fb_tree}
\end{figure*}

\subsection{Voraz: Matriz de Descarte Secuencial y Método Piloto}
El algoritmo Voraz no toma decisiones miopes; usa \textit{Lookahead} y la heurística WMDD. En cada turno proyecta el impacto de cada decisión sobre los tablones restantes.

En el \textbf{Turno 1 ($t=0$)}, el algoritmo evalúa los 5 tablones. Si elige $T_0$ primero (Costo Actual = 0), el costo simulado futuro es mayor a 80. Si elige $T_2$ primero (Costo Actual = 16), el costo futuro simulado baja a 28. Al comparar todas las proyecciones (ver \textbf{Figura \ref{fig:voraz_timeline}}), elige $T_2$ porque minimiza el daño total proyectado ($16 + 28 = 44$). En el Turno 2 repite la lógica con los 4 tablones restantes y elige $T_3$, y así hasta completar la secuencia.

\begin{figure*}[htbp]
\centering
\begin{tikzpicture}[
  scale=0.8, transform shape,
  node distance=2cm,
  box/.style={rectangle, draw=blue!80, fill=blue!5, thick, rounded corners, minimum width=3.2cm, minimum height=1.2cm, align=center},
  arrow/.style={->, thick, >=stealth}
]
  % Nodos de la línea de tiempo principal
  \node[box] (t0) {Turno 1 ($t=0$)\\ Evaluar los 5 Tablones};
  \node[box, right of=t0, xshift=4.5cm] (t1) {Turno 2 ($t=5$)\\ $T_2$ ya fue regado};
  \node[box, right of=t1, xshift=4.5cm] (t2) {Turno 3 ($t=9$)\\ $T_3$ ya fue regado};
  \node[right of=t2, xshift=2.8cm, font=\bfseries] (end) {\dots Termina en 44};

  % Flechas principales
  \draw[arrow] (t0) -- node[above, font=\footnotesize] {Elige definitivamente $T_2$} (t1);
  \draw[arrow] (t1) -- node[above, font=\footnotesize] {Elige definitivamente $T_3$} (t2);
  \draw[arrow] (t2) -- (end);

  % Cajas de simulación (arriba)
  \node[rectangle, draw=red!50, fill=red!5, dashed, align=left, above of=t0, yshift=1.5cm, font=\scriptsize] (sim0) {
    \textbf{Simulaciones WMDD (Score Total = Actual + Futuro):}\\
    $\times$ Simular $T_0 \to 0 + \text{Futuro}(84) = 84$\\
    $\times$ Simular $T_1 \to 8 + \text{Futuro}(56) = 64$\\
    \textcolor{green!60!black}{\textbf{\checkmark Simular $T_2 \to 16 + \text{Futuro}(28) = \mathbf{44}$}}
  };
  \draw[->, dashed, red!50] (t0) -- (sim0);

  \node[rectangle, draw=red!50, fill=red!5, dashed, align=left, above of=t1, yshift=1.5cm, font=\scriptsize] (sim1) {
    \textbf{Simulaciones WMDD (Score Total):}\\
    $\times$ Simular $T_0 \to 18 + \text{Futuro}(46) = 64$\\
    $\times$ Simular $T_4 \to 24 + \text{Futuro}(28) = 52$\\
    \textcolor{green!60!black}{\textbf{\checkmark Simular $T_3 \to 16 + \text{Futuro}(28) = \mathbf{44}$}}
  };
  \draw[->, dashed, red!50] (t1) -- (sim1);

\end{tikzpicture}
\caption{Línea de tiempo del algoritmo Voraz. El simulador WMDD calcula una proyección de penalidades y elige el camino que minimice el score total (costo local + costo futuro estimado).}
\label{fig:voraz_timeline}
\end{figure*}

\subsection{Programación Dinámica: Matriz de Ensamblaje Bottom-Up}
El diccionario de estados de la PD almacena subproblemas ya resueltos. A diferencia de la Fuerza Bruta, aquí la memoria es clave. La \textbf{Figura \ref{fig:pd_graph}} muestra que cada entrada del diccionario responde a una sola pregunta: \textit{''¿Cuál es el costo mínimo para regar exactamente este subconjunto de tablones?''}.

La resolución se construye de abajo hacia arriba: el algoritmo calcula que regar $T_0$ al final tiene un costo de 24. Con eso en memoria, el estado $\{T_0, T_4\}$ simplemente suma el costo local de $T_4$ (4) al costo guardado de $\{T_0\}$ (24), obteniendo 28. Esta cadena continúa hasta el estado inicial con los 5 tablones, que confirma el óptimo global de 44. Si otra ruta necesita el costo de $\{T_0, T_4\}$, lo lee directamente del diccionario en $O(1)$, sin recalcular nada.

\begin{figure*}[htbp]
\centering
\begin{tikzpicture}[
  scale=0.8, transform shape,
  level distance=2.2cm,
  level 1/.style={sibling distance=5.5cm},
  level 2/.style={sibling distance=3cm},
  every node/.style={rectangle, draw=purple!80, fill=purple!5, thick, rounded corners, align=center, scale=0.85},
  edge from parent/.style={draw, thick, ->, >=stealth}
]
  \node {\textbf{Estado Inicial} (Llave en RAM)\\ $\{T_0, T_1, T_2, T_3, T_4\}$\\ $\implies$ \textbf{Guarda: Costo 44, Decisión $T_2$}}
    child {node {Sub-Llave: $\{T_0, T_1, T_3, T_4\}$\\ $\implies$ \textbf{Guarda: Costo 28, Decisión $T_3$}}
      child {node {Sub-Llave: $\{T_0, T_1, T_4\}$\\ $\implies$ \textbf{Guarda: Costo 28, Decisión $T_1$}}
          child {node {Sub-Llave: $\{T_0, T_4\}$\\ $\implies$ \textbf{Guarda: Costo 28, Decisión $T_4$}}
             child {node {Sub-Llave: $\{T_0\}$\\ $\implies$ \textbf{Costo 24}}
             }
          }
      }
      child {node[draw=gray!50, fill=gray!10] {Otras combinaciones\\ cacheadas $\dots$}}
    }
    child {node[draw=gray!50, fill=gray!10] {Otras sub-ramas cacheadas\\ $\{T_1, T_2, T_3, T_4\} \dots$}};
\end{tikzpicture}
\caption{Grafo de dependencias de subproblemas (Programación Dinámica). Cada caja es una entrada en el diccionario de memoización. El algoritmo construye el costo óptimo global a partir de subproblemas resueltos previamente.}
\label{fig:pd_graph}
\end{figure*}

\section{Validación Empírica}
Los tres algoritmos se ejecutaron sobre 19 casos de prueba y produjeron costos equivalentes a la salida esperada oficial. El Test 10 (con 8 tablones) merece atención particular: la salida esperada es 522, que la Fuerza Bruta y la PD confirman de forma exacta. El Voraz obtiene 532, lo que confirma su naturaleza de estrategia aproximada.

\section{Algoritmo 1 — Fuerza Bruta}
\subsection{Idea Conceptual}
La Fuerza Bruta construye todo el árbol de combinaciones posibles. En cada nivel de recursión elige un tablón no regado, calcula el costo de regarlo en el tiempo actual y baja recursivamente para resolver el resto de la finca.

\begin{algorithm}
\caption{Fuerza Bruta: Encontrar Óptimo}
\begin{algorithmic}[1]
\STATE \textbf{procedure} ExplorarSinMemo($pendientes, tiempo$)
    \IF{$pendientes = \emptyset$}
        \RETURN $0, []$
    \ENDIF
    \STATE $mejor\_costo \gets \infty$, $mejor\_orden \gets []$
    \FOR{$idx \in pendientes$}
        \STATE $costo\_local \gets \text{CostoTablon}(idx, tiempo)$
        \STATE $resto \gets pendientes \setminus \{idx\}$
        \STATE $t\_nuevo \gets tiempo + tr_{idx}$
        \STATE $costo\_futuro, orden \gets \text{ExplorarSinMemo}(resto, t\_nuevo)$
        \STATE $total \gets costo\_local + costo\_futuro$
        \STATE $orden\_candidata \gets [idx] \cup orden$
        \IF{$total < mejor\_costo$}
            \STATE $mejor\_costo \gets total$
            \STATE $mejor\_orden \gets orden\_candidata$
        \ELSIF{$total = mejor\_costo \text{ y } orden\_candidata > mejor\_orden$}
            \STATE $mejor\_costo \gets total$
            \STATE $mejor\_orden \gets orden\_candidata$
        \ENDIF
    \ENDFOR
    \RETURN $mejor\_costo, mejor\_orden$
\STATE \textbf{end procedure}
\end{algorithmic}
\end{algorithm}

\subsection{Complejidad Computacional (Fuerza Bruta)}
Para $n$ tablones, el algoritmo evalúa todas las permutaciones posibles.

\begin{itemize}
    \item \textbf{Complejidad Temporal ($O(n \cdot n!)$):} Sea $T(k)$ el tiempo para explorar un nodo con $k$ tablones pendientes. El algoritmo itera sobre los $k$ elementos y en cada iteración construye la estructura \texttt{resto} en $O(k)$ y llama recursivamente a $T(k-1)$:
    \begin{equation}
    T(k) = k \cdot \Big( T(k-1) + O(k) \Big)
    \end{equation}
    Desplegando la recurrencia:
    \begin{align*}
    T(n) &= n \cdot T(n-1) + n \cdot O(n) \\
         &= n(n-1)T(n-2) + n(n-1)O(n-1) + n \cdot O(n) \\
         &= \sum_{i=0}^{n-1} \frac{n!}{(n-i)!} \cdot O(n-i)
    \end{align*}
    El número total de nodos en el árbol de recursión está acotado por $\sum_{k=0}^n \frac{n!}{k!} \approx e \cdot n!$. Como cada nodo realiza $O(n)$ operaciones en el peor caso:
    \begin{equation}
    T(n) = O(n \cdot n!)
    \end{equation}
    
    \item \textbf{Complejidad Espacial ($O(n^2)$):} La ejecución es una búsqueda en profundidad (DFS) con profundidad máxima $h = n$. En cada nivel $i$ se instancian arreglos temporales de tamaño $O(i)$. La memoria acumulada en cualquier instante es:
    \begin{equation}
    S(n) = \sum_{i=1}^{n} O(i) = O\left( \frac{n(n+1)}{2} \right) = O(n^2)
    \end{equation}
\end{itemize}

\subsection{Consistencia y Desempate Lexicográfico}
Un aspecto importante de la implementación es el manejo de empates: cuando dos permutaciones producen el mismo costo mínimo, los tres algoritmos deben elegir el mismo camino para poder verificar sus resultados entre sí. Para esto implementamos un \textbf{desempate lexicográfico descendente}: si dos órdenes cuestan lo mismo, se elige el de índices mayores. Por ejemplo, si $[2, 1, 0]$ y $[0, 1, 2]$ tienen el mismo costo, el algoritmo siempre elige $[2, 1, 0]$ porque empieza con 2, que es mayor que 0. Esto garantiza consistencia entre los tres algoritmos y simplifica la verificación de correctitud.

\subsection{Demostración Matemática (Fuerza Bruta)}

La correctitud de la Fuerza Bruta se sigue directamente de que explora la totalidad del espacio de soluciones. Formalmente:

\textbf{Teorema 1.} \textit{El algoritmo de Fuerza Bruta garantiza encontrar el óptimo global $\Pi^*$.}

\begin{proof}
Sea $S_n$ el espacio de permutaciones con $|S_n| = n!$, y $f: S_n \to \mathbb{R}^+$ la función de costo.
El algoritmo genera el conjunto $E$ de todas las evaluaciones realizadas. Por definición recursiva, $\forall \Pi \in S_n \implies \Pi \in E$, por lo tanto $E \equiv S_n$.
El costo óptimo se define como $C^* = \min_{\Pi \in S_n} f(\Pi)$, y el algoritmo retorna $\hat{C} = \min_{\Pi \in E} f(\Pi)$.
Como $E \equiv S_n$:
\begin{equation}
\hat{C} = \min_{\Pi \in E} f(\Pi) \equiv \min_{\Pi \in S_n} f(\Pi) \implies \hat{C} = C^*
\end{equation}
\end{proof}

\section{Algoritmo 2 — Voraz}
\subsection{Heurística WMDD y Método Piloto}
El algoritmo Voraz evita la miopía clásica de estas técnicas usando un método de simulación \textit{Lookahead} combinado con la heurística WMDD (Weighted Modified Due Date) \cite{morton1993heuristic}.

En lugar de elegir el tablón más barato en el momento actual, el algoritmo simula el futuro: por cada candidato, finge haberlo regado ya y evalúa cómo quedarían los demás tablones. Para ordenar ese ''resto'' rápidamente, usa la fórmula WMDD, que toma el tiempo de vida restante de cada tablón y lo divide por su prioridad económica. Si a un tablón de alta prioridad le queda poco tiempo, ese cociente es pequeño, lo que indica urgencia. El algoritmo suma el costo inmediato más el costo proyectado de la simulación y elige el candidato con el menor total. Así evita el error clásico de optimizar el paso actual a costa de penalizar los siguientes.

\begin{algorithm}
\caption{Algoritmo Voraz (Lookahead con WMDD)}
\begin{algorithmic}[1]
\STATE $orden \gets []$, $tiempo \gets 0$, $pendientes \gets \{0..n-1\}$
\WHILE{$pendientes \neq \emptyset$}
    \STATE $mejor\_candidato \gets \text{NULO}$, $min\_proyeccion \gets \infty$
    \FOR{$c \in pendientes$}
        \STATE $costo\_local \gets \text{CostoTablon}(c, tiempo)$
        \STATE $t\_futuro \gets tiempo + tr_c$
        \STATE $resto \gets pendientes \setminus \{c\}$
        \STATE $resto\_ordenado \gets \text{WMDD}(resto, t\_futuro)$
        \STATE $costo\_futuro \gets \text{Simular}(resto\_ordenado, t\_futuro)$
        \STATE $puntaje \gets costo\_local + costo\_futuro$
        \IF{$puntaje < min\_proyeccion$}
            \STATE $min\_proyeccion \gets puntaje$
            \STATE $mejor\_candidato \gets c$
        \ELSIF{$puntaje = min\_proyeccion \text{ y } c > mejor\_candidato$}
            \STATE $mejor\_candidato \gets c$ 
        \ENDIF
    \ENDFOR
    \STATE $orden.\text{push}(mejor\_candidato)$
    \STATE $tiempo \gets tiempo + tr_{mejor\_candidato}$
    \STATE $pendientes \gets pendientes \setminus \{mejor\_candidato\}$
\ENDWHILE
\RETURN $orden$
\end{algorithmic}
\end{algorithm}


\subsection{Complejidad Computacional (Voraz con WMDD)}
El lookahead introduce una anidación de simulaciones que aleja al algoritmo de las métricas de un voraz clásico.

\begin{itemize}
    \item \textbf{Complejidad Temporal ($O(n^3 \log n)$):} En el turno con $k$ tablones pendientes, el algoritmo evalúa $k$ candidatos. Para cada candidato aplica \texttt{timsort} sobre $k-1$ elementos ($O(k \log k)$) y luego recorre el resultado linealmente ($O(k)$). El trabajo por turno es:
    \begin{equation}
    W(k) = k \cdot \Big[ O(k \log k) + O(k) \Big] \approx O(k^2 \log k)
    \end{equation}
    El tiempo total suma el trabajo de todos los turnos:
    \begin{equation}
    T(n) = \sum_{k=1}^{n} O(k^2 \log k)
    \end{equation}
    Acotando $\log k \leq \log n$ como constante y usando la identidad de la suma de cuadrados:
    \begin{equation}
    T(n) \le \log n \left( \sum_{k=1}^{n} k^2 \right) = \log n \cdot \frac{n(n+1)(2n+1)}{6}
    \end{equation}
    \begin{equation}
    T(n) = O(n^3 \log n)
    \end{equation}

    \item \textbf{Complejidad Espacial ($O(n)$):} El algoritmo mantiene solo colecciones de tamaño $\leq n$ en todo momento. No hay almacenamiento exponencial ni recursión profunda:
    \begin{equation}
    E(n) = O(n) + O(n) + O(n-1) = O(n)
    \end{equation}
    Esto le permite procesar fincas grandes sin consumo excesivo de memoria.
\end{itemize}

\subsection{Limitaciones Teóricas}
Al tomar decisiones locales sin explorar todas las permutaciones, el Voraz no garantiza el óptimo global. En casos como el Test 10, el error es de aproximadamente $1.9\%$. Sin embargo, produce resultados muy cercanos al óptimo sin volverse inviable computacionalmente.

\subsection{Demostración Matemática (Algoritmo Voraz)}

En cada turno el Voraz toma la mejor decisión disponible según su estimación del futuro, pero esa estimación puede diferir del valor real. Esto es lo que formalmente lo hace subóptimo a nivel global:

\textbf{Teorema 2.} \textit{El algoritmo Voraz garantiza el óptimo local en cada paso $t$, pero puede ser globalmente subóptimo.}

\begin{proof}
Sea $P_t$ el conjunto de tablones disponibles en el instante $t$. Sea $f_{loc}: P_t \times \mathbb{R} \to \mathbb{R}$ el costo unitario, y $f_{sim}: \mathcal{P}(P_t) \times \mathbb{R} \to \mathbb{R}$ la función heurística WMDD de proyección.
El algoritmo elige $c^*$ mediante:
\begin{equation}
c^* = \arg\min_{\forall c \in P_t} \Big[ f_{loc}(c, t) + f_{sim}(P_t \setminus \{c\}, t + tr_c) \Big]
\end{equation}
Esta selección garantiza el óptimo condicionado al estado actual. Sin embargo, dado el carácter heurístico de WMDD, existe $P_t \subset F$ tal que $f_{sim}(P_t, t) \neq f_{opt}(P_t, t)$, donde $f_{opt}$ es el costo real futuro.
Acumulando el error sobre los $n$ pasos:
\begin{equation}
C_V \ge \sum_{k=1}^n f_{loc}(\pi_k^*, t_k) \implies C_V \ge C^*
\end{equation}
\end{proof}

\section{Algoritmo 3 — Programación Dinámica}
\subsection{Subestructura Óptima y Memoización}
La Programación Dinámica (PD) supera a la Fuerza Bruta almacenando resultados de subproblemas ya resueltos. El estado de la finca depende únicamente de \textit{qué tablones quedan por regar}, sin importar el orden histórico en que se regaron los anteriores.

Sea $C(P, t)$ el costo óptimo para los tablones pendientes en $P$ al tiempo $t$:
\begin{equation}
C(P, t) = \min_{i \in P} \left[ \text{Costo}(i, t) + C(P \setminus \{i\}, t + tr_i) \right]
\end{equation}

Si el algoritmo llega a un conjunto $P$ que ya fue resuelto antes, simplemente devuelve su valor del diccionario en $O(1)$.

La diferencia con la Fuerza Bruta es que la PD no recalcula subproblemas. Si ya sabemos cuál es el costo mínimo para regar exactamente los tablones 3 y 4 a partir del día 10, ese resultado queda guardado. La próxima vez que por cualquier ruta lleguemos al día 10 con los tablones 3 y 4 pendientes, no se ejecuta ningún cálculo extra: se lee directamente del diccionario. Esto elimina millones de cálculos redundantes.

\begin{algorithm}
\caption{Programación Dinámica Top-Down}
\begin{algorithmic}[1]
\STATE \textbf{Global:} Diccionarios $memo\_valor$ y $memo\_decision$
\STATE \textbf{procedure} CalcularValorOptimo($pendientes, tiempo$)
    \IF{$pendientes = \emptyset$}
        \RETURN $0$
    \ENDIF
    \IF{$pendientes \in memo\_valor$}
        \RETURN $memo\_valor[pendientes]$
    \ENDIF
    \STATE $mejor\_costo \gets \infty$, $mejor\_idx \gets -1$
    \FOR{$idx \in pendientes$}
        \STATE $costo\_local \gets \text{CostoTablon}(idx, tiempo)$
        \STATE $resto \gets pendientes \setminus \{idx\}$
        \STATE $costo\_resto \gets \text{CalcularValorOptimo}(resto, tiempo + tr_{idx})$
        \STATE $total \gets costo\_local + costo\_resto$
        \IF{$total < mejor\_costo$}
            \STATE $mejor\_costo \gets total$
            \STATE $mejor\_idx \gets idx$
        \ELSIF{$total = mejor\_costo \text{ y } idx > mejor\_idx$}
            \STATE $mejor\_idx \gets idx$
        \ENDIF
    \ENDFOR
    \STATE $memo\_valor[pendientes] \gets mejor\_costo$
    \STATE $memo\_decision[pendientes] \gets mejor\_idx$
    \RETURN $mejor\_costo$
\STATE \textbf{end procedure}
\end{algorithmic}
\end{algorithm}


\subsection{Complejidad Computacional (Programación Dinámica)}
La memoización reduce el crecimiento factorial a uno exponencial, al evitar que los mismos subproblemas se calculen más de una vez.

\begin{itemize}
    \item \textbf{Complejidad Temporal ($O(n \cdot 2^n)$):} Para $n$ tablones existen $2^n$ subconjuntos posibles. Gracias al diccionario con acceso $O(1)$, ningún subconjunto se evalúa dos veces. Cuando el algoritmo procesa un subconjunto de tamaño $k$, itera $k$ veces para aplicar el operador $\min$. El tiempo total es:
    \begin{equation}
    T(n) = \sum_{k=1}^{n} \binom{n}{k} \cdot O(k)
    \end{equation}
    Aplicando la identidad combinatoria $k\binom{n}{k} = n\binom{n-1}{k-1}$:
    \begin{align*}
    T(n) &= O\left( \sum_{k=1}^{n} n \binom{n-1}{k-1} \right) \\
         &= O\left( n \sum_{j=0}^{n-1} \binom{n-1}{j} \right)
    \end{align*}
    Por el Teorema del Binomio de Newton $(1+1)^{n-1} = 2^{n-1}$:
    \begin{equation}
    T(n) = O\left(n \cdot 2^{n-1}\right) \equiv O(n \cdot 2^n)
    \end{equation}
    
    \item \textbf{Complejidad Espacial ($O(n \cdot 2^n)$):} Los diccionarios deben almacenar los resultados de los $2^n$ subproblemas posibles. Dado que cada entrada tiene un tamaño proporcional a $k$, la memoria total requerida es:
    \begin{equation}
    S(n) = \sum_{k=1}^{n} \binom{n}{k} \cdot O(k) = O(n \cdot 2^n)
    \end{equation}
    Este crecimiento exponencial en memoria explica por qué el algoritmo agota la RAM al superar $n \approx 25$: $2^{25}$ entradas son aproximadamente 33 millones de registros en el diccionario.
\end{itemize}

\subsection{Demostración Matemática (Programación Dinámica)}

La correctitud de la PD se basa en el principio de optimalidad de Bellman: si la solución a un subproblema es óptima, entonces la solución global que la incluye también lo es. Por inducción fuerte sobre el tamaño del conjunto pendiente, se puede demostrar que cada entrada del diccionario contiene el costo mínimo real para ese subconjunto.

\textbf{Teorema 3.} \textit{La Programación Dinámica garantiza el óptimo global $C^*$.}

\begin{proof}
Procedemos por inducción fuerte sobre $k = |P|$.

\textbf{Base ($k=0$):} Si $P = \emptyset$, no hay tablones pendientes y $C(\emptyset, t) = 0$, que es trivialmente óptimo.

\textbf{Hipótesis inductiva ($k-1$):} Para todo $P'$ con $|P'| = k-1$, asumimos que $C(P', t') = C^*_{opt}(P', t')$.

\textbf{Paso inductivo ($k$):} Sea $|P| = k$. Por la ecuación de Bellman:
\begin{equation}
C(P, t) = \min_{i \in P} \Big[ f(i, t) + C(P \setminus \{i\}, t + tr_i) \Big]
\end{equation}
Para todo $i \in P$, el conjunto $P \setminus \{i\}$ tiene tamaño $k-1$. Por la hipótesis inductiva, $C(P \setminus \{i\}, \cdot) = C^*_{opt}$. Como $f(i,t)$ es un costo determinista exacto y el operador $\min$ se evalúa sobre todos los $i \in P$:
\begin{equation}
C(P, t) = \min_{\forall i \in P} \Big[ f(i, t) + C_{opt}^*(P \setminus \{i\}, t + tr_i) \Big] \equiv C^*
\end{equation}
Por inducción fuerte, la propiedad se cumple para todo $k \leq n$.
\end{proof}

\section{Diseño de Casos de Prueba}
Los 19 casos de prueba cubren escenarios variados:
\begin{itemize}
    \item Casos pequeños ($n=5$) donde el déficit hídrico es bajo y fácil de manejar.
    \item Escenarios con prioridades mixtas, donde elegir primero tablones rápidos perjudica a tablones de alta prioridad ($p_i=4$).
    \item Distribuciones complejas con empates teóricos, como el Test 10, resueltos por la regla lexicográfica.
    \item Instancias de escala mayor ($n=20$ y $n=25$) para evaluar los límites de cada algoritmo.
\end{itemize}
La estructura valida la exactitud de la Fuerza Bruta y la PD, y permite medir el margen de error admisible del Voraz.

\section{Comparación Global de los Algoritmos}
La Programación Dinámica produce resultados exactos hasta $n \approx 18$ (230 ms para $n=15$). A partir de $n \geq 20$ el consumo de memoria se vuelve prohibitivo. El Voraz resuelve el caso $n=25$ en 3 milisegundos con variaciones menores al 3\%. La Fuerza Bruta, aunque exacta, no es viable en la práctica para $n \geq 12$.

\section{Detalles de Implementación}
La arquitectura del sistema separa claramente las responsabilidades. Un \textbf{Frontend} moderno gestiona la interfaz, y un \textbf{Backend} en Node.js expone endpoints REST que reciben las fincas en formato JSON. El motor algorítmico está codificado en \textbf{Python} e invocado mediante pipes desde el backend. Los módulos están organizados por función: \texttt{modelo\_riego.py} contiene los tres algoritmos, \texttt{costo\_tablon.py} implementa la función de penalidad y \texttt{utils\_estado.py} maneja la representación del estado de tablones pendientes.

\section{Análisis de Correctitud}
El principio de optimalidad de Bellman \cite{bellman1966dynamic} fundamenta la correctitud de la PD: un camino óptimo está compuesto por sub-caminos que también son óptimos. Al visitar cada subconjunto de tablones de forma top-down con memoización, cada resultado guardado es globalmente correcto. Por su parte, la Fuerza Bruta garantiza correctitud por exhaustividad: evalúa el 100\% de las permutaciones, por lo que no puede dejar el óptimo sin evaluar.

\section{Resultados Experimentales}
\begin{table}[htbp]
\centering
\caption{Costos y Rendimiento (Muestra Aleatoria de Pruebas)}
\resizebox{\columnwidth}{!}{%
\begin{tabular}{@{}cccccc@{}}
\toprule
\textbf{Test} & \textbf{$n$} & \textbf{Costo Ideal} & \textbf{Costo V} & \textbf{T. FB (ms)} & \textbf{T. PD (ms)} \\ \midrule
1 & 5 & 44 & 44 & 0 & 0 \\
8 & 7 & 294 & 294 & 12 & 0 \\
10 & 8 & 522 & 532 & 92 & 0 \\
12 & 10 & 526 & 542 & 8214 & 4 \\
15 & 15 & 1030 & 1037 & Timeout & 230 \\
19 & 25 & 3594 & 3782 & Timeout & Timeout \\ \bottomrule
\end{tabular}%
}
\end{table}

A medida que $n$ aumenta, el Voraz presenta un margen de error creciente pero contenido. El tiempo de la Fuerza Bruta crece un factor de $\sim n$ por cada tablón adicional (de 92 ms a 8 segundos al pasar de $n=8$ a $n=10$). El uso de memoria de la PD se convierte en el factor limitante en los últimos tests.

\subsection{Análisis Detallado del Comportamiento por Algoritmo}

\textbf{Entradas pequeñas (Test 1 y Test 8):}
Para $n=5$ y $n=7$, los tres algoritmos producen el mismo resultado (44 y 294). Con 120 y 5.040 combinaciones respectivamente, la Fuerza Bruta y la PD evalúan todo el espacio en tiempo despreciable. El Voraz también acierta porque con horizontes cortos su estimación WMDD coincide con el óptimo real.

\textbf{El punto de divergencia (Test 10, $n=8$):}
La Fuerza Bruta y la PD coinciden en \textbf{522}, evaluando las 40.320 combinaciones posibles. El Voraz obtiene \textbf{532}, fallando por 10 puntos. La causa es la \textit{miopía heurística}: en un turno intermedio, el Voraz evaluó que postergar un tablón $A$ de alta prioridad causaría un alto daño futuro, y lo regó de inmediato. Sin embargo, existía una combinación posterior que permitía acomodar $A$ sin penalidades. La PD y la Fuerza Bruta encontraron esa combinación porque exploraron el espacio completo; el Voraz no, porque solo trabajó con una estimación.

\textbf{El límite de escala (Test 15 y Test 19):}
Para $n=15$, la Fuerza Bruta supera el tiempo límite: $1.3$ billones de permutaciones exceden la capacidad de cómputo disponible. La PD resuelve el mismo caso en 230 ms encontrando el óptimo \textbf{1030}. En $n=25$ la PD también falla, esta vez por memoria: $2^{25}$ subconjuntos generan cerca de 33 millones de entradas en el diccionario. El Voraz resuelve este caso en milisegundos con un costo de \textbf{3782}, lo que lo convierte en la única opción viable para fincas grandes.

\section{Discusión Académica}
En un contexto agroindustrial real, no siempre vale la pena buscar el óptimo exacto si calcularlo requiere tiempos de cómputo inviables. El Voraz, basado en heurísticas con respaldo teórico del área de Scheduling, produce soluciones con un error de entre 1\% y 3\% que se calculan en milisegundos, incluso en hardware de gama baja. Este balance entre calidad de solución y tiempo de respuesta lo hace más adecuado para escenarios de producción reales que los algoritmos exactos.

\section{Trabajos Futuros}
Para instancias con $n > 20$ donde se requiera mayor exactitud, una dirección natural es explorar metaheurísticas como Algoritmos Genéticos o Simulated Annealing. También se identifican dos optimizaciones concretas sobre los algoritmos actuales:
\begin{itemize}
    \item \textbf{Poda (Branch and Bound):} En la Fuerza Bruta, mantener el mejor costo parcial encontrado y descartar ramas que ya lo superan, reduciendo el número de evaluaciones sin sacrificar exactitud.
    \item \textbf{Paralelización del lookahead:} En el Voraz, el ciclo de simulaciones de candidatos puede ejecutarse en paralelo con procesamiento multi-core, reduciendo el tiempo total en fincas grandes.
\end{itemize}

\section{Conclusiones}
El problema de secuenciamiento de riego con penalizaciones por tiempo no tiene solución exacta en tiempo polinomial dado su carácter NP-Hard. La Programación Dinámica reduce la complejidad de $O(n!)$ a $O(n \cdot 2^n)$, lo que la hace viable hasta $n \approx 18$. Para fincas más grandes, el Voraz con lookahead predictivo resulta ser la alternativa más práctica: escala bien, consume poca memoria y produce resultados con márgenes de error aceptables para el contexto agroindustrial.

\section{Referencias}
\begin{thebibliography}{00}

\bibitem{cormen2009introduction}
T. H. Cormen, C. E. Leiserson, R. L. Rivest, y C. Stein, \textit{Introduction to Algorithms}, 3ra ed. Cambridge, MA: MIT Press, 2009. \\
\url{https://mitpress.mit.edu/9780262033848/introduction-to-algorithms/}

\bibitem{graham1979optimization}
R. L. Graham, E. L. Lawler, J. K. Lenstra, y A. H. G. Rinnooy Kan, ``Optimization and Approximation in Deterministic Sequencing and Scheduling: A Survey,'' \textit{Annals of Discrete Mathematics}, vol. 5, pp. 287--326, 1979. \\
DOI: \href{https://doi.org/10.1016/S0167-5060(08)70356-X}{10.1016/S0167-5060(08)70356-X}

\bibitem{pinedo2016scheduling}
M. L. Pinedo, \textit{Scheduling: Theory, Algorithms, and Systems}, 5ta ed. Nueva York, NY: Springer, 2016. \\
DOI: \href{https://doi.org/10.1007/978-3-319-26580-3}{10.1007/978-3-319-26580-3}

\bibitem{bellman1966dynamic}
R. Bellman, ``Dynamic Programming,'' \textit{Science}, vol. 153, no. 3731, pp. 34--37, 1966. \\
DOI: \href{https://doi.org/10.1126/science.153.3731.34}{10.1126/science.153.3731.34}

\bibitem{morton1993heuristic}
T. E. Morton y D. W. Pentico, \textit{Heuristic Scheduling Systems}. Nueva York, NY: Wiley, 1993. \\
DOI: \href{https://doi.org/10.1002/9781118627372}{10.1002/9781118627372}

\end{thebibliography}

\end{document}